\section{Related Work}
\label{sec:related_work}

Our study lies at the intersection of three research strands: maintenance and
service-contract optimization, reinforcement learning for sequential
decision-making in industrial operations, and deep policy-gradient methods.


\paragraph{Maintenance and service-contract optimization.}
Service-contract design and maintenance scheduling have a long tradition in
operations research, where they are typically formulated as integer programs,
renewal-theoretic models or stochastic dynamic programs
\cite{jardine2006managing,dekker1996applications}. These approaches provide
principled solutions under strong modelling assumptions but scale poorly to
high-dimensional portfolios with stochastic failures and evolving contractual
structures, and typically produce static rather than adaptive policies. In
industrial practice, simple threshold rules and segment-based heuristics
dominate deployments \cite{gits1992structuring}, which motivates the search
for data-driven alternatives that can capture the temporal feedback between
coverage decisions and downstream costs.


\paragraph{Reinforcement learning for industrial operations.}
Reinforcement learning has emerged as a natural framework for sequential
decision-making in industrial and financial operations, thanks to its
ability to learn from interactions without requiring an explicit model of
the environment \cite{sutton2018reinforcement}. Recent applications include
risk-sensitive trading agents trained with language-model-augmented rewards
\cite{benhenda2025}, deep RL for efficient liquidity provisioning in
decentralized finance \cite{xu2025}, and dynamic pricing and inventory
management \cite{boute2022deep}. Complementary efforts benchmark large
language models as decision-making agents in financial and business
prediction tasks~\cite{benhenda2026lookaheadbench,%
chen2025vcbenchbenchmarkingllmsventure,benhenda2026ycbench,zeng2025,%
chandak2026}. These works illustrate both the promise of learned
decision-making under uncertainty and the recurring difficulties of reward
specification, sample efficiency and safe deployment---difficulties that
are also central to our setting.


\paragraph{Reinforcement learning in healthcare operations.}
Within healthcare, RL has been applied to treatment recommendation, dynamic
treatment regimes and operational decision-making such as ICU ventilator
weaning and resource allocation \cite{yu2021reinforcement}. Most of these
applications emphasize safety, interpretability and offline learning from
historical logs, all constraints that also apply when managing a fleet of
medical devices. To the best of our knowledge, contract-coverage
optimization for medical equipment has not yet been cast as an RL problem in
the published literature, which motivates the present work.


\paragraph{Policy-gradient methods.}
At the algorithmic level, our approach builds on deep actor-critic methods.
Proximal Policy Optimization \cite{schulman2017ppo} provides an empirically
stable policy-gradient procedure through a clipped surrogate objective,
which we couple with Generalized Advantage Estimation
\cite{schulman2016gae}. Running observation normalization and advantage
normalization, which we also use, are standard components of modern PPO
implementations \cite{andrychowicz2021matters}. Our contribution is not
algorithmic but applicative: we design an MDP formulation, a simulation
environment and an evaluation protocol tailored to contractual-coverage
optimization, and we situate the learned policy against a random baseline
and a strong rule-based heuristic.
